% !TeX program = lualatex
% !TeX encoding = utf8
% !TeX spellcheck = uk_UA
% !TeX root = ../main.tex

%=========================================================
\Opensolutionfile{answer}[\currfilebase/\currfilebase-Answers]
\Writetofile{answer}{\protect\section*{\nameref*{\currfilebase}}}
\chapter{Магнітне поле в речовині}\label{\currfilebase}\hypertarget{MagFieldMedia}{}
\graphicspath{{\currfilebase/Pictures/}}
%=========================================================
\epigraph{\Alexander Сама Земля~--- великий магніт.}{Вільям Гільберт, \enquote{Про магніт}
	(De Magnete, 1600)}


%%% --------------------------------------------------------
\section*{Гіпотеза Ампера}
%%% --------------------------------------------------------


Ми\index{Ампера гіпотеза} вже знаємо, що магнітне поле діє на рухомі заряджені частинки та на рамки зі струмом.
Але \emph{чому воно діє також і на шматок звичайного заліза} --- скажімо, на цвях, який
сам по собі не є магнітом? Відповідь на це запитання дав ще А.~Ампер~\cite{Ampere1820}: він припустив, що
\emph{всередині магніту теж течуть струми}.

Якщо взяти магнітну стрілку чи магніт у руки, жодних струмів ми, звісно, не відчуваємо:
отже, ці струми циркулюють \emph{усередині} речовини і ніколи не виходять назовні. Що це
за струми, Ампер, звісно, не знав --- адже про будову атома тоді не було відомо нічого.
Сучасній науці вже відомо, що звичайна речовина складається з атомів: усередині атомів є
позитивно заряджені ядра, навколо яких рухаються негативно заряджені електрони, причому
самі електрони, до того ж --- маленькі магнітні \enquote{стрілочки} (мають власний
магнітний момент, спін). Отже, рух електронів усередині атомів --- це і є ті самі
\emph{струми}, про існування яких здогадався Ампер; магнітне поле діє на ці струми, а
отже --- і на речовину загалом (рис.~\ref{pic:AmpereHypotesis}).

%=========================================================
\begin{figure}[h!]\centering
%---------------------------------------------------------
\begin{subfigure}[t]{0.45\linewidth}\centering
    \localinput{unmagnetized.tikz}
    \caption{Ненамагнічена речовина (хаотичний розподіл молекулярних струмів).}
    \label{pic:unmagnetized}
\end{subfigure}
\hfill % Додає відступ між картинками
%---------------------------------------------------------
\begin{subfigure}[t]{0.45\linewidth}\centering
    \localinput{magnetized.tikz}
    \caption{Намагнічена речовина (впорядкований розподіл молекулярних струмів).}
    \label{pic:magnetized}
\end{subfigure}
%---------------------------------------------------------
\caption{Гіпотеза Ампера: магнетизм речовини зумовлений мікроскопічними
    (молекулярними) струмами всередині неї:
    \protect\subref{pic:unmagnetized}~--- ненамагнічена речовина;
    \protect\subref{pic:magnetized}~--- намагнічена речовина.}
\label{pic:AmpereHypotesis}
\end{figure}
%=========================================================


%% --------------------------------------------------------
\section{Основні поняття}\hypertarget{magnetics-basic-concepts}{}
%% --------------------------------------------------------


%%% --------------------------------------------------------
\subsection*{Мікрополе та середнє поле}
%%% --------------------------------------------------------


У речовині магнітне поле формується як зовнішнім полем, так і струмами, що циркулюють у самій речовині. На мікрорівні (тобто на відстанях порядку
розміру атома й менше) це поле різко змінюється в часі й просторі --- таке поле називають \emphx{мікрополем} $\Bfield_\text{micro}$. Однак якщо
провести усереднення за \emph{фізично нескінченно малим} об'ємом $\Delta V$ (малим порівняно з масштабом, на якому поле помітно змінюється
макроскопічно, але великим порівняно з розмірами й відстанями між окремими атомами --- тобто таким, що вміщує величезну кількість атомів), отримаємо
\emphx{середнє поле}:
\begin{equation*}
		\langle\Bfield\rangle = \frac1{\Delta V}\iiint\limits_{\Delta V} \Bfield_\text{micro}\, dV.
\end{equation*}
Середнє поле змінюється значно повільніше за мікрополе, оскільки усереднення за $\Delta V$ згладжує швидкі просторові й часові коливання мікрополя. Надалі, як і
раніше з полем $\Efield$ в діелектрику, під $\Bfield$ ми завжди розумітимемо саме це \emphx{середнє поле}.

Цей прийом той самий, що й для поля $\Efield$ в діелектриках (розділ~\ref{FieldInDielectrics}): у макроскопічній електродинаміці ми ніколи не працюємо з полем окремого атома, а лише з полем, усередненим за фізично нескінченно малим об'ємом.

\subsection*{Струми провідності та молекулярні струми}

Усі струми, створювані рухомими зарядами в речовині, поділяють на дві принципово різні групи:
\begin{enumerate}
	\item \emphx{Струми провідності} --- пов'язані з перенесенням \emphx{вільних} носіїв заряду (тих самих, що ми вивчали в розділі~\ref{SteadyCurrent}); це \emphx{макроскопічні} струми, які течуть крізь речовину; їх задають або вимірюють ззовні (джерело ЕРС, амперметр).
	\item \emphx{Молекулярні струми}\index{Молекулярні струми} --- зумовлені орбітальним рухом електронів в атомах (молекулах) та їхнім спіном (власним моментом імпульсу), а
	      також, значно меншою мірою, магнітними моментами ядер. Це \emph{внутрішні}, властиві самій речовині струми.
\end{enumerate}

%% --------------------------------------------------------
\section{Намагнічування}\hypertarget{magnetization}{}
%% --------------------------------------------------------

\subsection*{Вектор намагніченості}

\emphx{Вектором намагніченості}\index{Намагніченість} (\emphx{намагніченістю}) називають величину, що характеризує сумарний магнітний момент одиниці об'єму речовини:
\begin{equation}\label{eq:MagnetizationDef}
	\tcbhighmath{
		\vect J = \frac1{\Delta V}\sum\limits_{i\,\in\,\Delta V} \vect p_{m,i},
	}
	\qquad \vect p_{m,i} \text{ --- момент $i$-го атома,}
\end{equation}
де $\Delta V$ --- той самий \emph{фізично нескінченно малий} об'єм.
Намагніченість називають \emphx{однорідною}, якщо вектор $\vect J$ не залежить від точки в речовині; якщо $\vect J \neq \const$ --- \emphx{неоднорідною}.

Вектор $\vect J$ --- магнітний аналог вектора поляризації $\vect P$ (розділ~\ref{FieldInDielectrics}): там $\vect P$ --- густина електричних дипольних моментів, тут $\vect J$ --- густина магнітних. Далі логіка розділу (зв'язані струми, допоміжний вектор $\Hfield$, граничні умови) повторює логіку розділу про діелектрики; точну відповідність величин зведено в таблицю~\ref{tab:DielMagAnalogy}.

\subsection*{Основна задача теорії магнітостатики в речовині}

Пряме обчислення поля в намагніченій речовині вимагало б підсумовування магнітних полів від колосальної кількості (порядку числа Авогадро)
молекулярних струмів --- завдання абсолютно нездійсненне практично. Основна ідея теорії магнітостатики в речовині --- \emph{вилучити} молекулярні
струми з рівнянь, замінивши їх іншою, макроскопічно \emph{вимірюваною} величиною --- вектором намагніченості $\vect J$. Це, за аналогією з тим, що
ми вже робили для зв'язаних зарядів у діелектриках, істотно спрощує розрахунок полів у магнетиках.

\subsection*{Зв'язок намагніченості з молекулярними струмами}

Виділимо в намагніченій речовині косокутний циліндр таких малих розмірів, що намагніченість $\vect J$ у ньому можна вважати однорідною, і орієнтуємо його
так, щоб $\vect J$ була напрямлена по нормалі $\vect n$ до торців (рис.~\ref{pic:Cylinder}); тобто молекулярні струми течуть у площинах, паралельних торцям. У товщі циліндра кільцеві молекулярні струми сусідніх атомів течуть у протилежних напрямках і взаємно компенсуються --- в
усіх точках, \emph{окрім} точок бічної поверхні (рис.~\ref{pic:CylinderTop}, вигляд торця): там компенсація неповна, і залишається нескомпенсований
поверхневий струм, що тече бічною поверхнею циліндра.

\begin{figure}[h!]
	\centering
	\begin{minipage}[t]{0.48\linewidth}\centering
		\localinput{cylinder.tikz}
		\caption{Косокутний циліндр у намагніченій речовині}\label{pic:Cylinder}
	\end{minipage}\hfill
	\begin{minipage}[t]{0.48\linewidth}\centering
		\localinput{cylindertop.tikz}
		\caption{Вигляд торця циліндра: кільцеві струми компенсуються всюди, крім бічної поверхні}\label{pic:CylinderTop}
	\end{minipage}
\end{figure}

Магнітний момент такого циліндрика, з одного боку, дорівнює $\vect p_m = \vect J V$ (за означенням намагніченості), а з іншого --- визначається
нескомпенсованим поверхневим струмом $I_m$, що тече по бічній поверхні, за формулою для магнітного моменту витка~\eqref{eq:MagMomentLoop}
(розділ~\ref{MagFieldVacuum}), $\vect p_m = \frac1c I_m S\vect n$, де $S$ --- площа торця циліндра, $\vect n$ --- нормаль до торця, а $\theta$ ---
кут між твірною $\vect\tau$ циліндра і цією нормаллю. Об'єм косокутного циліндра дорівнює $V = S\,\ell\cos\theta$, де $\ell$ --- довжина
твірної. Прирівнюючи обидва вирази й проєктуючи рівність на напрям твірної:
\begin{equation*}
	\vect J\, V = \frac1c I_m S\vect n \quad\Rightarrow\quad (\vect J\cdot\vect\tau)\, S\,\ell\cos\theta = \frac1c I_m\, S\,(\vect n\cdot\vect\tau),
	\qquad \vect n\cdot\vect\tau = \cos\theta.
\end{equation*}
Скорочуючи на $S\cos\theta$ і ділячи на $\ell$, знаходимо \emphx{поверхневу густину молекулярного струму} $i_m = I_m/\ell$ (струм на одиницю довжини твірної):
\begin{equation}\label{eq:MolecularSurfaceCurrent}
	\tcbhighmath{
		i_m = c\, \vect J\cdot\vect\tau.
	}
\end{equation}

Тут $i_m$~--- струм через одиницю довжини твірної; у векторній формі $\vect i_m = c\,[\vect J\times\vect\nu]$, де $\vect\nu$~--- зовнішня нормаль до поверхні речовини. Звідси видно, що \emph{поверхневий} молекулярний струм перпендикулярний до намагніченості, $\vect i_m\perp\vect J$, і визначається \emph{тангенціальною} складовою $\vect J$. Це відрізняється від електричного аналога: зв'язаний поверхневий заряд $\sigma'=P_n$ визначається \emph{нормальною} складовою $\vect P$ (розділ~\ref{FieldInDielectrics}). Причина в тому, що $\vect P$ і заряд пов'язані через дивергенцію (скалярне джерело), а $\vect J$ і струм~--- через ротор.

%% --------------------------------------------------------
\section{Теорема про циркуляцію в речовині}\hypertarget{circulation-theorem-in-matter}{}
%% --------------------------------------------------------

\subsection*{Циркуляція вектора намагніченості}

Виберемо тепер у речовині \emph{довільний} замкнений контур $L$ (рис.~\ref{pic:Jloop}). Молекулярний струм, що перетинає елемент контуру $d\symbfscr{l}$,~--- це якраз нескомпенсований струм на бічній поверхні тонкого циліндра з твірною $d\symbfscr{l}$, тож за формулою~\eqref{eq:MolecularSurfaceCurrent} він дорівнює $dI_m = i_m\, d\ell = c\,\vect J\cdot d\symbfscr{l}$, тому повний молекулярний струм, що
перетинає (пронизує) контур,
\begin{equation}\label{eq:CirculationJ}
	I_m = c\oint\limits_L \vect J\cdot d\symbfscr{l}.
\end{equation}

\begin{SCfigure}[2][h!]\centering
	\localinput{Jloop.tikz}
	\caption{Довільний замкнений контур $L$ у намагніченій речовині\label{pic:Jloop}}
\end{SCfigure}

За теоремою Стокса цей вираз перетворюється так:
\begin{equation*}
	I_m = c\oint\limits_L \vect J\cdot d\symbfscr{l} = c\iint\limits_S \rot\vect J\cdot d\vect S,
\end{equation*}
де $S$ --- будь-яка поверхня, натягнута на контур $L$. З іншого боку, той самий струм $I_m$ можна виразити через його \emph{об'ємну} густину
$\vect j_m$: $I_m = \iint_S \vect j_m\cdot d\vect S$. Порівнюючи два вирази, отримуємо \emph{зв'язок молекулярного струму з вектором
	намагніченості в диференціальній формі}:
\begin{equation}\label{eq:MolecularCurrentDensity}
	\tcbhighmath{
		\vect j_m = c\rot\vect J.
	}
\end{equation}

\subsection*{Вектор напруженості магнітного поля}

Циркуляцію повного магнітного поля $\Bfield$ у речовині породжують \emph{усі} струми --- і струми провідності $I$, і молекулярні струми $I_m$:
\begin{equation*}
	\oint\limits_L \Bfield\cdot d\symbfscr{l} = \frac{4\pi}c\left(I + I_m\right) = \frac{4\pi}c\left(I + c\iint\limits_S \rot\vect J\cdot d\vect
	S\right).
\end{equation*}
Щоб \emph{вилучити} невимірюваний молекулярний струм $I_m$ з цього рівняння, введемо допоміжну величину --- \emphx{напруженість магнітного поля}\index{Напруженість магнітного поля}:
\begin{equation}\label{eq:HDef}
	\tcbhighmath{
		\Hfield = \Bfield - 4\pi\vect J.
	}
\end{equation}
Підставляючи це означення в попереднє рівняння, молекулярний струм скорочується, і \emphx{теорема про циркуляцію в речовині}\index{Теорема про циркуляцію!у магнетику} набуває вигляду
\begin{equation}\label{eq:CirculationHInt}
	\tcbhighmath{
		\oint\limits_L \Hfield\cdot d\symbfscr{l} = \frac{4\pi}c I.
	}
\end{equation}
Еквівалентна диференціальна форма цієї теореми:
\begin{equation}\label{eq:CirculationHDiff}
	\tcbhighmath{
		\rot\Hfield = \frac{4\pi}c\vect j.
	}
\end{equation}

У гауссовій системі $\Bfield$ і $\Hfield$ мають однакову розмірність: одиниця $\Bfield$~--- гаус (Гс), одиниця $\Hfield$~--- ерстед (Е); у вакуумі $1~\text{Е}=1~\text{Гс}$.

\begin{Attention}[unbreakable]
	Вектор $\Hfield$~--- \emph{допоміжна} величина, що слугує лише для спрощення рівнянь; принципово важливо, що \emph{його циркуляція визначається тільки струмами провідності}~--- тими, які ми задаємо чи вимірюємо ззовні, а не невідомими молекулярними струмами. За роллю в рівняннях $\Hfield$ подібний до вектора $\Dfield$: той дав змогу записати теорему Гаусса лише через вільні заряди (розділ~\ref{FieldInDielectrics}). Але за формою означення відповідність інша (таблиця~\ref{tab:DielMagAnalogy}).
\end{Attention}

%% --------------------------------------------------------
\section{Магнітна сприйнятливість речовини}\hypertarget{magnetic-susceptibility}{}
%% --------------------------------------------------------

%% --------------------------------------------------------
\subsection*{Лінійні ізотропні середовища}
%% --------------------------------------------------------

Розглянемо \emphx{лінійні ізотропні} середовища --- ті, в яких
намагніченість пропорційна полю, а магнітні властивості однакові в усіх
напрямках. Середовище (речовина), що реагує на зовнішнє магнітне поле, називатимемо
\emphx{магнетиком}. Якщо поле не надто сильне, а в середовищі немає
початкової (залишкової) намагніченості, намагніченість, як показує дослід,
прямо пропорційна полю:
\begin{equation}\label{eq:LinearMedium}
	\vect J = \chi_m\Hfield,
\end{equation}
де $\chi_m$ --- \emphx{магнітна сприйнятливість}\index{Магнітна сприйнятливість} середовища (безрозмірна
величина). Підставляючи це в означення~\eqref{eq:HDef},
$\Bfield = \Hfield + 4\pi\vect J = (1+4\pi\chi_m)\Hfield$, тобто
\begin{equation}\label{eq:BMuH}
	\tcbhighmath{
		\Bfield = \mu\Hfield, \qquad \mu = 1+4\pi\chi_m,
	}
\end{equation}
де величину $\mu$ називають \emphx{магнітною проникністю середовища}\index{Магнітна проникність}.

Тепер можна точно зіставити електричну й магнітну теорії (таблиця~\ref{tab:DielMagAnalogy}). За \emph{формою} означень відповідність така: $\Dfield=\Efield+4\pi\vect P$ \,$\leftrightarrow$\, $\Bfield=\Hfield+4\pi\vect J$, тобто $\Dfield\leftrightarrow\Bfield$, $\Efield\leftrightarrow\Hfield$, $\vect P\leftrightarrow\vect J$, $\epsilon\leftrightarrow\mu$; саме тому однаково виглядають граничні умови й закони заломлення. За \emph{роллю} в теоремі про вільні джерела $\Hfield$ подібний до $\Dfield$ (обидва означено так, щоб молекулярні джерела випали), але зв'язані джерела виражаються по-різному: через дивергенцію для $\vect P$ і через ротор для $\vect J$.

\begin{table}[h!]\small\centering
	\caption{Відповідність величин у діелектриках і магнетиках (гауссова система)\label{tab:DielMagAnalogy}}
	\begin{tblr}{
			colspec={X[l, m]X[c,m]X[c,m]},
			row{odd} = {gray!10},
			row{1} = {c, m, fg=white, bg=themecolorlight, font=\bfseries},
			hlines,
			vlines
		}
		                                      & Діелектрик                                                            & Магнетик                                                               \\
		Густина дипольних моментів            & $\vect P$                                                             & $\vect J$                                                              \\
		Зв'язані джерела                      & $\rho'=-\divg\vect P$                                                & $\vect j_m=c\rot\vect J$                                              \\
		Допоміжний вектор                     & $\Dfield=\Efield+4\pi\vect P$                                        & $\Bfield=\Hfield+4\pi\vect J$                                         \\
		Лінійне середовище                    & $\vect P=\chi\Efield$, $\Dfield=\epsilon\Efield$, $\epsilon=1+4\pi\chi$ & $\vect J=\chi_m\Hfield$, $\Bfield=\mu\Hfield$, $\mu=1+4\pi\chi_m$ \\
		Теорема з вільними джерелами          & $\divg\Dfield=4\pi\rho$                                              & $\rot\Hfield=\dfrac{4\pi}c\vect j$                                    \\
		Неперервні на межі без вільних джерел & $E_\tau$, $D_n$                                                       & $H_\tau$, $B_n$                                                        \\
		Закон заломлення                      & $\dfrac{\tg\alpha_2}{\tg\alpha_1}=\dfrac{\epsilon_2}{\epsilon_1}$    & $\dfrac{\tg\alpha_2}{\tg\alpha_1}=\dfrac{\mu_2}{\mu_1}$             \\
	\end{tblr}
\end{table}

%% --------------------------------------------------------
\subsection*{Класифікація магнетиків}
%% --------------------------------------------------------

Магнетиком є \emph{будь-яка} речовина. Термін \enquote{діелектрик} виділяє
лише частину речовин і протиставляється провіднику; термін \enquote{магнетик}
не протиставляється нічому~--- магнітно інертних речовин не існує. Причина~--- діамагнітний відгук: зовнішнє поле змінює рух електронів у \emph{кожному} атомі й наводить момент, напрямлений проти поля (його механізм розглянуто нижче, у~розділі про ларморову прецесію).

Наявність мікрострумів, однак, ще не означає наявності \emphx{магнітного
моменту}: у більшості атомів орбітальні й спінові моменти електронів в
основному стані взаємно компенсуються, і сумарний момент атома дорівнює нулю
--- точний аналог неполярних молекул у діелектриках. Тому речовини
відрізняються не наявністю мікрострумів, а тим, чи скомпенсовані їхні моменти
за відсутності поля і як ці струми відгукуються на його ввімкнення.

Залежно від знака й величини магнітної сприйнятливості виділяють такі
основні класи лінійних магнетиків:
\begin{enumerate}
	\item Якщо $\chi_m > 0$ (а отже, $\mu>1$), магнетик називають
	      \emphx{парамагнетиком}\index{Парамагнетики}. Парамагнітні властивості мають, наприклад,
	      \ce{Al}, \ce{Pt}, \ce{FeCl2}, \ce{O2}, більшість лужних і лужноземельних металів.
	      (У~металах парамагнетизм здебільшого зумовлений електронами
	      провідності~--- це так званий парамагнетизм Паулі, слабкий і~майже
	      незалежний від температури; закон Кюрі, виведений нижче, описує речовини
	      з~власними магнітними моментами атомів чи молекул, як \ce{O2} чи солі
	      заліза.)
	\item Якщо $\chi_m < 0$ (а отже, $\mu<1$), магнетик називають
	      \emphx{діамагнетиком}\index{Діамагнетики}. Діамагнетиками є \ce{Bi}, \ce{Cu}, \ce{Au}, \ce{Sb}, \ce{Si},
	      \ce{H2O}, \ce{H2}, \ce{N2} тощо.
\end{enumerate}

За порядком величини магнітна сприйнятливість пара- і діамагнетиків мала: $|\chi_m|\sim10^{-7}$--$10^{-3}$,
тобто $\mu$ відрізняється від одиниці лише на частки відсотка й менше --- зовсім інший масштаб, ніж $\mu\gg1$
у феромагнетиків, про які йдеться нижче.

І парамагнетики, і діамагнетики належать до класу лінійних ізотропних
середовищ, розглянутих вище. Граничний випадок діамагнетика~--- надпровідник, у якому $\chi_m=-1/4\pi$ і $\mu=0$: поле повністю виштовхується з його об'єму.

Деякі речовини здатні зберігати намагніченість $\vect J$ навіть за
\emph{відсутності} зовнішнього поля --- для них лінійне
співвідношення~\eqref{eq:LinearMedium} узагалі не виконується. Такі
речовини називають \emphx{феромагнетиками}\index{Феромагнетики} (\ce{Fe}, \ce{Co}, \ce{Ni} та
багато їхніх сплавів). Для них відношення $\mu=B/H$ залежить від $H$ і від передісторії зразка; характерні значення $\mu$~--- тисячі й навіть десятки тисяч, $\mu\gg1$. Детальніше про феромагнетики~--- наприкінці цього розділу.




%% --------------------------------------------------------
\section{Граничні умови}\hypertarget{magnetic-boundary-conditions}{}
%% --------------------------------------------------------

\subsection*{Перша гранична умова (для нормальної складової $\Bfield$)}

Застосуємо\index{Граничні умови!на межі магнетиків} теорему Гаусса для вектора $\Bfield$ (яка, нагадаємо, залишається справедливою в будь-якому середовищі, оскільки джерел магнітного поля,
на відміну від електричного, --- магнітних зарядів --- не існує) до нескінченно малого циліндра, що охоплює невелику ділянку межі розділу
двох магнетиків (рис.~\ref{pic:BC1Mag}). Вважаючи $d\vect S_2 = -d\vect S_1 = \vect n\, dS$ ($\vect n$ --- нормаль від середовища 1 до
середовища 2, як на рисунку):
\begin{equation*}
	\oiint\limits_S \Bfield\cdot d\vect S = 0 \quad\Rightarrow\quad \Bfield_1\cdot d\vect S_1 + \Bfield_2\cdot d\vect S_2 = 0.
\end{equation*}

\begin{figure}[h!]\centering
	\localinput{bc1.tikz}
	\caption{Циліндр Гаусса на межі розділу двох магнетиків\label{pic:BC1Mag}}
\end{figure}

Звідси перша гранична умова:
\begin{equation}\label{eq:BC1Magnetic}
	\tcbhighmath{
		B_{1n} = B_{2n}
	}
\end{equation}
--- \emph{нормальна складова вектора $\Bfield$ не зазнає стрибка} при переході через межу розділу середовищ (дослівний магнітний аналог
неперервності $D_n$ за відсутності вільного заряду на межі, розділ~\ref{FieldInDielectrics}).

\subsection*{Друга гранична умова (для тангенціальної складової $\Hfield$)}

Застосуємо теорему про циркуляцію~\eqref{eq:CirculationHInt} до нескінченно малого прямокутного контуру $ABCD$, довгі сторони якого проходять
на нескінченно малій відстані по обидва боки від межі розділу (рис.~\ref{pic:BC2Mag}). Поверхневий струм провідності (якщо він там є)
характеризуємо одиничним вектором $\vect b$ уздовж напряму струму й лінійною густиною $i_b$; контур орієнтуємо \emph{впоперек} струму, так щоб
струм пронизував його, тобто вздовж орта $\vect\tau = \vect b\times\vect n$. Обходячи контур так, що $d\vect r_2 = \vect\tau\, d\ell$ (сторона в
середовищі 2) і $d\vect r_1 = -d\vect r_2$:
\begin{equation*}
	\oint\limits_L \Hfield\cdot d\symbfscr{l} = \frac{4\pi}c I \quad\Rightarrow\quad \Hfield_1\cdot d\vect r_1 + \Hfield_2\cdot d\vect r_2 =
	\frac{4\pi}c\, i_b\, d\ell.
\end{equation*}

\begin{figure}[h!]\centering
	\localinput{bc2.tikz}
	\caption{Контур циркуляції на межі розділу двох магнетиків\label{pic:BC2Mag}}
\end{figure}

Звідси друга гранична умова:
\begin{equation}\label{eq:BC2MagneticScalar}
	\tcbhighmath{
		H_{2\tau} - H_{1\tau} = \frac{4\pi}c\, i_b.
	}
\end{equation}
Оскільки $(\Hfield_2-\Hfield_1)\cdot(\vect b\times\vect n)=\vect b\cdot[\vect n\times(\Hfield_2-\Hfield_1)]$, а контур іншої орієнтації струм не пронизує (циркуляція по ньому дорівнює нулю), умову записують у векторному вигляді:
\begin{equation}\label{eq:BC2MagneticVector}
	\tcbhighmath{
		\vect n\times\left(\Hfield_2 - \Hfield_1\right) = \frac{4\pi}c\, \vect i_b.
	}
\end{equation}

\begin{Attention}[unbreakable]
	Тангенціальна складова $\Hfield$ зазнає стрибка \emph{тільки} там, де по поверхні розділу течуть \emphx{струми провідності} $i_b$ (наприклад, обмотка, намотана безпосередньо на межі двох магнетиків). Без таких струмів $H_{1\tau}=H_{2\tau}$, а $B_\tau=\mu H_\tau$ узагалі кажучи \emph{стрибає}, бо $\mu_1\neq\mu_2$~--- так само, як на межі діелектриків без вільного заряду неперервна $E_\tau$, а $D_\tau=\epsilon E_\tau$ стрибає (таблиця~\ref{tab:DielMagAnalogy}).
\end{Attention}

\subsection*{Заломлення силових ліній на межі магнетиків}

Розглянемо межу розділу двох магнетиків з проникностями $\mu_1$ і $\mu_2$, на якій немає струмів провідності (рис.~\ref{pic:RefractionM}). Нехай
силова лінія утворює кут $\alpha_1$ з нормаллю в середовищі 1 і кут $\alpha_2$ --- в середовищі 2. Граничні умови дають:
\begin{equation*}
	\begin{cases}
		H_{1\tau}=H_{2\tau} \ \Rightarrow\ H_1\sin\alpha_1 = H_2\sin\alpha_2, \\
		B_{1n}=B_{2n} \ \Rightarrow\ \mu_1H_1\cos\alpha_1 = \mu_2H_2\cos\alpha_2.
	\end{cases}
\end{equation*}

\begin{figure}[h!]\centering
	\localinput{refraction.tikz}
	\caption{Заломлення силових ліній $\Hfield$ і $\Bfield$ на межі розділу двох
		магнетиків\label{pic:RefractionM}}
\end{figure}

Ділячи ці рівняння одне на одне, отримуємо \emphx{закон заломлення} силових ліній, однаковий і для $\Bfield$, і для $\Hfield$:
\begin{equation}\label{eq:RefractionLaw}
	\tcbhighmath{
		\frac{\tg\alpha_2}{\tg\alpha_1} = \frac{\mu_2}{\mu_1}.
	}
\end{equation}
У середовищі з більшим значенням $\mu$ силові лінії відхиляються на більший кут від нормалі до межі розділу.

\subsection*{Магнітне екранування}

На\index{Магнітне екранування} заломленні силових ліній ґрунтується явище \emph{магнітного захисту (екранування)}. Якщо замкнену оболонку з феромагнетика (наприклад,
залізну) внести в зовнішнє магнітне поле, то, оскільки $\mu_\text{залізо} \gg \mu_\text{повітря}\approx 1$, за законом
заломлення~\eqref{eq:RefractionLaw} силові лінії всередині оболонки відхиляються майже до дотичної площини до її поверхні --- тобто
\emph{\enquote{втягуються} в товщу оболонки} і концентруються переважно в ній самій, майже не проникаючи у внутрішню порожнину
(рис.~\ref{pic:MagShield}). У результаті поле всередині порожнини виявляється істотно ослабленим порівняно із зовнішнім~--- але не
нульовим: на відміну від електростатичного захисту, магнітна оболонка лише \emph{послаблює} поле, тим сильніше, чим більші $\mu$ і~товщина
стінки (за порядком величини~--- у~$\sim\mu d/R$ разів, де $d$~--- товщина, $R$~--- радіус оболонки; це так, якщо $\mu d/R\gg1$). У змінних полях поле додатково послаблюють вихрові струми в провідній оболонці.

\begin{SCfigure}[1][h!]\centering
	\includegraphics[width=0.45\linewidth]{MagShield}
	\caption{Магнітне екранування: силові лінії концентруються в товщі
		феромагнітної оболонки\label{pic:MagShield}}
\end{SCfigure}

\begin{Attention}[unbreakable]
	Магнітне екранування --- принципово \emph{інше} явище, ніж електростатичне:
	провідник у електростатиці \emph{виштовхує} поле зі свого об'єму, тоді як
	феромагнітна оболонка, навпаки, \emph{затягує} лінії поля в себе, немов
	пропонуючи їм \enquote{шлях найменшого опору}. Тому магнітний захист
	вимагає \emph{суцільної замкненої} оболонки з високопроникного матеріалу
	--- на відміну від електростатичного, для якого досить навіть тонкої
	провідної сітки. Цей принцип застосовують для захисту чутливих приладів
	(магнітометрів, детекторів) від зовнішніх завад.
\end{Attention}



%% --------------------------------------------------------
\section{Мікроскопічна природа діа-, пара- і феромагнетизму}\hypertarget{dia-para-ferromagnetism-microscopic}{}
%% --------------------------------------------------------

Усе, що було в цьому розділі до цього моменту, --- феноменологія: $\vect J$, $\Hfield$, $\chi_m$ введені через
дослідні факти й формальні означення, без жодного припущення про будову атома.
Тепер з'ясуємо, звідки вона береться ---
тобто розкриємо мікроскопічну природу діа- і парамагнетизму, а наостанок з'ясуємо, звідки береться значно сильніше й принципово інше явище ---
спонтанне впорядкування у феромагнетиках.



%% --------------------------------------------------------
\subsection*{Природа діамагнетизму: прецесія Лармора}
%% --------------------------------------------------------



Скористаємося класичною (борівською) моделлю атома, де електрон рухається коловою орбітою навколо ядра --- саме в цій моделі діамагнетизм найпростіше побачити кількісно.\footnote{Це лише наочна оцінка. Строго
кажучи, класична статистична фізика дає для системи зарядів у~магнітному полі \emph{нульову} намагніченість (теорема Бора--ван~Левен~\cite{Bohr1911,VanLeeuwen1921}):
діамагнетизм, як і~парамагнетизм та феромагнетизм,~--- квантове явище. Модель зі \enquote{збереженою} орбітою неявно спирається на
квантову стійкість атома й~дає правильний за порядком і~знаком результат.} За формулою~\eqref{eq:GyromagneticClassical} (розділ~\ref{MagFieldVacuum}) його орбітальний
магнітний момент пов'язаний з орбітальним моментом імпульсу $\vect L$ гіромагнітним відношенням $\vect p = \gamma\vect L$, $\gamma=-e/(2m_ec)$. У
зовнішньому полі $\Bfield$ на цей момент діє момент сили~\eqref{eq:TorqueDipoleAnalogy}, який змінює $\vect L$ за рівнянням
\begin{equation*}
	\frac{d\vect L}{dt} = \vect p\times\Bfield = \gamma\left(\vect L\times\Bfield\right) = \vect\omega_L\times\vect L, \qquad
	\vect\omega_L\equiv-\gamma\Bfield.
\end{equation*}
Це --- рівняння рівномірної \emphx{прецесії}\index{Лармора прецесія} вектора $\vect L$ (а разом з ним --- і всієї площини орбіти) навколо напрямку $\Bfield$ з кутовою
швидкістю
\begin{equation}\label{eq:LarmorFreqMag}
	\tcbhighmath{
		\omega_L = |\gamma|B = \frac{eB}{2m_ec}
	}
\end{equation}
--- \emphx{ларморовою частотою}, однаковою для будь-якого електрона незалежно від радіуса й швидкості його орбіти (\emphx{теорема Лармора}\index{Лармора теорема}). Це справедливо, поки $\omega_L$ мала порівняно з частотою обертання електрона в атомі, тобто для $B\ll10^{9}$~Гс~--- для лабораторних полів із великим запасом.

%---------------------------------------------------------
\begin{figure}[h!]\centering
    \localinput{larmor_precession.tikz}
\caption{Прецесія орбіти електрона в зовнішньому полі $\Bfield$: магнітний
момент $\vect{p}_m$ (а разом з ним і площина орбіти) обертається навколо
$\Bfield$ з ларморовою частотою $\omega_L$. Ця прецесія еквівалентна
додатковому коловому струму $I=-e\omega_L/(2\pi)$, який створює індукований
магнітний момент $\Delta\vect p$, напрямлений проти $\Bfield$.}
\label{pic:LarmorPrecession}
\end{figure}
%---------------------------------------------------------


До ввімкнення поля електрон рухався по коловій орбіті; щойно поле накладено,
ця орбіта починає прецесувати (рис.~\ref{pic:LarmorPrecession}).  Прецесія орбіти --- це додатковий круговий рух заряду $-e$ з частотою $\omega_L$ навколо осі, паралельної $\Bfield$, тобто додатковий коловий
струм $I=-e\omega_L/(2\pi)$ на відстані $\rho$ від цієї осі (якщо орбіта перпендикулярна до $\Bfield$, то $\rho$~--- радіус орбіти; у загальному випадку $\rho$ змінюється вздовж орбіти, і далі ми усереднимо $\rho^2$). За формулою для моменту колового витка~\eqref{eq:MagMomentLoop} цей струм створює
додатковий (\emphx{індукований}) магнітний момент
\begin{equation}\label{eq:InducedDiaMag}
	\Delta\vect{p}_m = \frac{I\pi\rho^2}{c}\,\frac{\Bfield}{B} = -\frac{e\omega_L\rho^2}{2c}\,\frac{\Bfield}{B} = -\frac{e^2\rho^2}{4m_ec^2}\, \Bfield,
\end{equation}
напрямлений, зверніть увагу, \emph{проти} поля $\Bfield$ --- незалежно від того, у який бік електрон обертався початково: це і є мікроскопічний
прояв правила Ленца.

\begin{Example}{Індукований момент із закону електромагнітної індукції*}\small
	Той самий результат~\eqref{eq:InducedDiaMag} випливає із закону
	електромагнітної індукції. При вмиканні поля потік крізь орбіту
	$\Phi=\pi r^2 B$ змінюється, і виникає ЕРС
	$\mathcal{E}=-(1/c)\,d\Phi/dt$. Вона створює дотичне поле
	$E_{\text{ін}}=r\dot B/(2c)$, яке за час $dt$ змінює швидкість електрона на
	\begin{equation*}
		dv = \frac{eE_{\text{ін}}}{m_e}\,dt
		= \frac{er}{2m_ec}\,dB .
	\end{equation*}
	Відповідна зміна струму $dI=-e\,dv/(2\pi r)$ дає
	\begin{equation*}
		d\Delta p_m = \frac{\pi r^2}{c}\,dI
		= -\frac{e^2 r^2}{4m_ec^2}\,dB
		\quad\Longrightarrow\quad
		\Delta p_m = -\frac{e^2 r^2}{4m_ec^2}\,B .
	\end{equation*}
	Знак мінус --- це правило Ленца: індукований момент напрямлений проти поля,
	яке його викликало. Тут орбіта перпендикулярна до $\Bfield$, тому $\rho=r$; зміна швидкості електрона однакова за напрямом для обох напрямів обертання, тож $\Delta p_m$ завжди напрямлений проти поля.
\end{Example}

Формула~\eqref{eq:InducedDiaMag} виведена для однієї класичної орбіти. У
квантовій механіці електрон не має траєкторії: його стан описується хвильовою
функцією, а заряд розподілений у просторі з густиною $-e|\psi|^2$. Тому
замість фіксованого $\rho$ слід підставити середнє $\langle\rho^2\rangle$ за
цією електронною хмарою. Електронна хмара сферично
симетрична, тож
\begin{equation*}
	\langle\rho^2\rangle = \tfrac23\langle r^2\rangle .
\end{equation*}
Підсумовуючи внесок усіх $Z$ електронів атома (нижче $\langle r^2\rangle$~--- середнє значення по них), отримуємо індукований момент атома,
а тоді й магнітну сприйнятливість речовини з концентрацією атомів $n$
(\emphx{формула Ланжевена для діамагнетизму}\index{Ланжевена формула для діамагнетизму}~\cite{Langevin1905}):
\begin{equation}\label{eq:LangevinDiaMag}
	\tcbhighmath{
		\chi_m^{\text{діа}} = -\frac{nZe^2\langle r^2\rangle}{6m_ec^2} < 0.
	}
\end{equation}

Ця величина не залежить від температури (розмір атома $\langle r^2\rangle$ від неї практично не залежить) і за порядком величини становить
$10^{-7}$--$10^{-6}$ (у гауссовій системі; у СІ відповідні значення в $4\pi$ разів більші) --- слабкий, але \emph{універсальний} ефект,
наявний в атомах будь-якої речовини, включно з пара- і феромагнетиками.

%% --------------------------------------------------------
\subsection*{Природа парамагнетизму: закон Кюрі}
%% --------------------------------------------------------

Якщо атом має власний (нескомпенсований) магнітний момент $p_0\neq0$, до щойно знайденого діамагнітного внеску додається значно потужніший
ефект --- часткова \emph{орієнтація} вже наявних моментів уздовж поля. Це --- дослівний магнітний аналог орієнтаційної поляризації полярних
діелектриків (розділ~\ref{FieldInDielectrics}, \S~\ref{sec:Langevin-Debye}): енергія моменту в полі дорівнює $U=-p_0B\cos\theta$
(формула~\eqref{eq:PotentialEnergyDipole}, розділ~\ref{MagFieldVacuum}), а рівноважний розподіл орієнтацій підпорядковується розподілу Больцмана.
У практично важливій межі слабкого поля й не надто низької температури ($p_0B\ll kT$) середня проекція моменту на напрямок поля лінійна за
полем, за тією самою викладкою, що й для електричного диполя~\eqref{eq:LangevinDebye}:
\begin{equation}\label{eq:ParaAvgMag}
	\langle p_z\rangle = \frac{p_0^2 B}{3kT}.
\end{equation}
Звідси намагніченість $n$ атомів в одиниці об'єму дорівнює $J=n\langle p_z\rangle$. Оскільки $\chi_m\ll1$, у~самій речовині $B\approx H$,
тож $\chi_m=J/H\approx J/B$, а сприйнятливість спадає обернено пропорційно температурі ---
\emphx{закон Кюрі}\index{Кюрі закон}~\cite{Curie1895} (пряма аналогія закону~\eqref{eq:CurieLaw} для полярних діелектриків):
\begin{equation}\label{eq:CurieLawMag}
	\tcbhighmath{
		\chi_m^{\text{пара}} = \frac{np_0^2}{3kT} \equiv \frac{C}T > 0,
	}
\end{equation}
де $C=np_0^2/3k$ --- \emphx{стала Кюрі}. Оскільки за кімнатної температури $\chi_m^{\text{пара}}$, як правило, на один-два порядки перевищує
$\bigl|\chi_m^{\text{діа}}\bigr|$, саме вона визначає знак і величину сумарної сприйнятливості речовини, щойно в її атомах з'являється власний момент, --- а
слабший діамагнітний внесок~\eqref{eq:LangevinDiaMag} лишається практично непомітним на його тлі. (У квантовій теорії проєкції магнітного моменту атома набувають лише дискретних значень, і замість $p_0^2$ стоїть величина того самого порядку: для атома зі спіном $\tfrac12$ --- $3\mu_B^2$, де $\mu_B=e\hbar/2m_ec$~--- магнетон Бора. Закон $\chi\propto1/T$ при цьому зберігається.)

\subsection*{Природа феромагнетизму: обмінна взаємодія і поле Вейса}

Закон Кюрі~\eqref{eq:CurieLawMag} описує \emph{незалежні} атомні моменти, що орієнтуються лише під дією зовнішнього поля й теплового хаосу. Але
у феромагнетику є локальні області --- \emphx{домени}\index{Домени!магнітні} --- типовим розміром від часток до сотень мікрометрів. У межах кожного домену магнітні
моменти всіх атомів спонтанно вишикувані в одному напрямку (рис.~\ref{pic:Domains}) --- без жодного зовнішнього поля. У різних доменах напрями
намагніченості, як правило, різні, тому сумарна намагніченість ненамагніченого зразка загалом близька до нуля. Доменна структура цілком аналогічна до доменної структури сегнетоелектриків (розділ~\ref{FieldInDielectrics}). Розбиття на домени вигідне енергетично: воно зменшує енергію магнітного поля, яке однорідно намагнічений зразок створював би в навколишньому просторі.

\begin{SCfigure}[1.5][h!]\centering
	\includegraphics[width=0.35\linewidth]{domains}
	\caption{Доменна структура феромагнетика.
    %: у кожному домені намагніченість стала, але напрями
	%	намагніченості різних доменів, як правило, різні
    \label{pic:Domains}}
\end{SCfigure}


Формальне пояснення спонтанної орієнтації магнітних моментів у межах домену запропонував П.~Вейс~\cite{Weiss1907}: припустимо, що на кожен атомний момент,
окрім зовнішнього поля $\Hfield$, діє ще й
\emphx{молекулярне поле}\index{Вейса молекулярне поле} $\Hfield_{\text{мол}} = \lambda\vect J$, пропорційне намагніченості \emph{сусідніх} атомів (тут $\lambda$ --- стала
Вейса). Підставивши це ефективне поле $\Hfield+\lambda\vect J$ у закон Кюрі~\eqref{eq:CurieLawMag} замість $\Hfield$ (для слабких полів, коли
$B\approx H$), дістаємо самоузгоджене рівняння для намагніченості:
\begin{equation*}
	J = \frac CT\left(H+\lambda J\right) \quad\Longrightarrow\quad \chi_m = \frac JH = \frac{C}{T-C\lambda}.
\end{equation*}
Позначивши $T_C\equiv C\lambda$ --- \emphx{температуру Кюрі} (для заліза $T_C\approx770\,^\circ$C), отримуємо \emphx{закон Кюрі--Вейса}\index{Кюрі--Вейса закон}:
\begin{equation}\label{eq:CurieWeissLaw}
	\tcbhighmath{
		\chi_m = \frac{C}{T-T_C}.
	}
\end{equation}

Вище $T_C$ речовина парамагнітна: власні моменти атомів орієнтовані
хаотично, і без зовнішнього поля намагніченість дорівнює нулю. Коли
$T\to T_C^+$, сприйнятливість~\eqref{eq:CurieWeissLaw} прямує до
нескінченності --- парамагнітний стан втрачає стійкість. Щойно $T$
опускається нижче $T_C$, будь-яке мале відхилення $\vect J$ від нуля
підсилюється молекулярним полем швидше, ніж його руйнує тепловий рух, ---
і в системі самочинно встановлюється ненульова намагніченість навіть без
зовнішнього поля. Це --- феромагнітний стан. (Зростання $J$ зупиняється, коли всі моменти вишикувано, тобто на насиченні; щоб описати це кількісно, лінійний закон Кюрі вже не годиться.)

Залишається питання: чим фізично зумовлене поле Вейса $\Hfield_{\text{мол}}$? Щоб упорядкування витримувало тепловий рух аж до кімнатної температури,
потрібне поле порядку $10^6$--$10^7$~Гс (з умови $p_0B\sim kT$). Натомість поле, яке реально створює сусідній атомний диполь на міжатомній
відстані, становить лише $\sim10^3$~Гс --- на три-чотири порядки менше. Отже, \emph{справжня} причина --- не магнітна взаємодія диполів, а
\emphx{обмінна взаємодія}\index{Обмінна взаємодія}~\cite{Heisenberg1928}. Це --- суто квантовий ефект, що не має класичного
аналога: енергія системи двох електронів залежить не лише від їхніх зарядів
і відстані, а й від взаємної орієнтації їхніх спінів, причому ця залежність
не зводиться до магнітної. Причина --- у принципі Паулі: повна хвильова
функція двох електронів мусить бути антисиметричною, тому за
\emph{паралельних} спінів антисиметричною мусить бути її координатна
частина. Це автоматично \emph{розводить} такі електрони в просторі й тим
самим \emph{зменшує} енергію їхнього кулонівського відштовхування
порівняно з випадком антипаралельних спінів.

Обмінна взаємодія \emph{короткодійна} (діє практично лише між найближчими сусідами в ґратці) і в різних речовинах може сприяти як \emph{паралельній}, так і \emph{антипаралельній} орієнтації сусідніх моментів. Якщо вона вирівнює їх паралельно~--- маємо \emphx{феромагнетик}; якщо антипаралельно, причому моменти рівні за величиною,~--- сумарна намагніченість зникає, і речовину називають \emphx{антиферомагнетиком}\index{Антиферомагнетики}; якщо ж антипаралельно, але моменти різні за величиною (наприклад, належать різним типам вузлів ґратки),~--- повної компенсації не відбувається, і виникає \emphx{феримагнетик}\index{Феримагнетики} (класичний приклад~--- магнетит \ce{Fe3O4}, магнітний залізняк, з~якого складалися перші відомі людству природні магніти). У всіх трьох випадках природа взаємодії та сама, квантова; різні лише \emph{знак і геометрія} її прояву.



%\emphx{Феромагнетиками} називають тверді речовини, здатні мати \emph{спонтанну намагніченість всередині окремих доменів} --- тобто намагнічені
%навіть за \emph{відсутності} зовнішнього поля. Типові представники --- залізо, кобальт, нікель та багато їхніх сплавів.

\subsection*{Гістерезис у феромагнетиках}

\emphx{Гістерезисом}\index{Гістерезис} називають неоднозначну петлеподібну залежність намагніченості $J$ феромагнетика від зовнішнього поля $H$ за його циклічної
зміни (рис.~\ref{pic:Hysteresis}).

\begin{SCfigure}[1.25][h!]\centering
	\localinput{hysteresis.tikz}
	\caption{Петля гістерезису феромагнетика: $J_r$ --- залишкова намагніченість, $H_c$ ---
		коерцитивна сила\label{pic:Hysteresis}}
\end{SCfigure}

Проходження петлі відбувається за такими стадіями:
\begin{enumerate}
	\item За достатньо великого поля $H$ намагніченість досягає \emphx{насичення}: усі домени вже зорієнтовані вздовж поля, і зростання $J$
	      практично припиняється, $J\to J_s=\const$. Індукція $\Bfield = \Hfield+4\pi\vect J$ при цьому й далі зростає --- але вже лише за
	      рахунок самого $\Hfield$, з таким самим нахилом, як у вакуумі.
	\item Коли поле зменшують до нуля ($H=0$), намагніченість \emph{не} повертається до нуля, а залишається на рівні залишкової намагніченості
	      $J_r$ --- саме тому постійні магніти взагалі можуть існувати.
	\item Щоб звести намагніченість до нуля, потрібно прикласти поле протилежного напрямку --- \emph{коерцитивну силу} $-H_c$.
	\item З подальшим зростанням поля в тому самому (від'ємному) напрямку намагніченість знову зростає за модулем до насичення $-J_s$.
	\item Коли від'ємне поле зменшують за модулем до нуля, намагніченість знову залишається ненульовою, на рівні $-J_r$, --- а при подальшому зростанні поля, через $+H_c$, до $+J_s$ цикл замикається.
\end{enumerate}

Причина петлі~--- у тому, як змінюється доменна структура. У слабких полях ростуть домени, намагнічені вздовж поля, за рахунок сусідніх: межі між доменами (доменні стінки) зміщуються. Цей рух \emph{необоротний}: стінки зачіплюються за дефекти кристала й долають їх стрибками, тому після зняття поля частина доменів лишається зорієнтованою.


Площа петлі має пряму фізичну інтерпретацію: енергія, що незворотно розсіюється у вигляді тепла в~одиниці об'єму за один цикл перемагнічування, дорівнює $w=\frac1{4\pi}\oint H\,dB=\oint H\,dJ$, тобто площі петлі $J(H)$ (вихрові струми дають окремий внесок). Тому вузька петля означає малі втрати, а широка~--- великі.

Залежно від ширини петлі гістерезису феромагнетики поділяють на дві категорії. \emphx{Магнітом'які}\index{Феромагнетики!магнітом'які і магнітотверді} матеріали (рис.~\ref{pic:Soft}), які застосовують як осердя в трансформаторах, електромагнітах та інших приладах, мають дуже
\emph{вузьку} петлю гістерезису. Вузькість петлі (мала коерцитивна сила $H_c$) означає малі втрати на перемагнічування --- важливо, наприклад,
для трансформаторів, які перемагнічуються десятки чи сотні разів за секунду.

\begin{figure}[h!]\centering
	%---------------------------------------------------------
	\captionbox{Вузька петля гістерезису магнітом'якого матеріалу\label{pic:Soft}}[0.48\linewidth]{
		\centering
		\includegraphics[height=7cm]{Soft}
	}
	\hfill
	%---------------------------------------------------------
	\captionbox{Широка петля гістерезису магнітотвердого матеріалу\label{pic:Hard}}[0.48\linewidth]{
		\centering
		\includegraphics[height=7cm]{Hard}
	}
	%---------------------------------------------------------
\end{figure}

Постійні ж магніти, навпаки, виготовляють із \emphx{магнітотвердих} матеріалів (рис.~\ref{pic:Hard}), наприклад із загартованої сталі, сплавів Nd--Fe--B, феритів, з
широкою петлею гістерезису й великою коерцитивною силою $H_c$ --- саме це забезпечує збереження намагніченості магніту з часом і під дією
випадкових зовнішніх завад.

%% --------------------------------------------------------
\section*{Підсумки}
\phantomsection
\addcontentsline{toc}{section}{Підсумки}
%% --------------------------------------------------------

Уся теорія магнітного поля в речовині побудована за єдиним лекалом, яке ми вже бачили в теорії діелектриків: (1)~ввести густину
мікроскопічних джерел поля (там --- $\vect P$, тут --- $\vect J$); (2)~виразити через неї зв'язані (молекулярні) джерела --- там через
дивергенцію ($\rho' = -\divg\vect P$), тут через ротор ($\vect j_m = c\rot\vect J$) --- відмінність, пов'язана з тим, що $\Efield$ потенціальне, а $\Bfield$ вихрове; (3)~ввести допоміжний вектор ($\Dfield$ або $\Hfield$), теорема для якого містить \emph{лише} вільні
(вимірювані) джерела; (4)~для лінійних середовищ пов'язати допоміжний і основний вектори простим коефіцієнтом ($\epsilon$ або $\mu$); точну відповідність величин дає таблиця~\ref{tab:DielMagAnalogy}. Розуміння
цієї загальної структури --- найкращий спосіб не заплутатися в позначеннях і знаках обох розділів курсу.

Мікроскопічно магнітні властивості речовини визначаються тим, чи мають її атоми власний магнітний момент. \emph{Діамагнетизм}~---
універсальний слабкий ефект: зовнішнє поле спричиняє ларморову прецесію електронних орбіт і~наводить момент, напрямлений проти поля
($\chi_m<0$, не залежить від температури). \emph{Парамагнетизм} виникає, коли атоми мають власні моменти, які поле частково орієнтує
всупереч тепловому руху (закон Кюрі, $\chi_m=C/T>0$). У~\emph{феромагнетиках} обмінна взаємодія~--- квантовий ефект, у~тисячі разів
сильніший за магнітну взаємодію диполів,~--- вишиковує моменти в~доменах самочинно; вище температури Кюрі речовина стає парамагнітною
(закон Кюрі--Вейса). Феромагнетики нелінійні й~мають гістерезис: площа петлі визначає втрати на перемагнічування, ширина~--- поділ на
магнітом'які й~магнітотверді матеріали. На межі магнетиків неперервні $B_n$ і~$H_\tau$ (без струмів провідності), а~на цьому
ґрунтується, зокрема, магнітне екранування феромагнітними оболонками.

%% --------------------------------------------------------
\Needspace{9\baselineskip}
\section*{Задачі}
\phantomsection
\addcontentsline{toc}{section}{Задачі}
%% --------------------------------------------------------

% задачі цього розділу короткі: не розриваємо їх між сторінками

\begin{problem}\label{prb:SolenoidParamagnet}
	\emph{Соленоїд з~магнетиком.} Знайдіть індукцію магнітного поля всередині нескінченного соленоїда з~густиною намотки $n$ витків на
	одиницю довжини, по витках якого тече струм $I$, якщо об'єм усередині соленоїда заповнено магнетиком із проникністю $\mu$. Яка частина
	поля створена молекулярними струмами?
	\begin{solution}
		$H=\dfrac{4\pi}{c}nI$ (теорема про циркуляцію для $\Hfield$), $B=\mu H=\dfrac{4\pi\mu}{c}nI$; молекулярні струми дають
		$4\pi J=(\mu-1)H$, тобто частку $(\mu-1)/\mu$ від $B$.
	\end{solution}
\end{problem}

\begin{problem}\label{prb:ParamagneticWire}
	\emph{Молекулярні струми в~парамагнітному провіднику.} Постійний струм $I$ тече вздовж довгого циліндричного провідника круглого
	перерізу радіуса $R$, виготовленого з~парамагнетика зі сприйнятливістю $\chi_m$; струм розподілений по перерізу рівномірно. Знайдіть
	об'ємний і~поверхневий молекулярні струми та з'ясуйте, як вони напрямлені один відносно одного й~відносно $I$.
	\begin{solution}
		За теоремою про циркуляцію $H=\dfrac{2Ir}{cR^2}$ (азимутально), $J=\chi_mH$. Об'ємний струм
		$j_{m,z}=\dfrac{c}{r}\dfrac{d(rJ)}{dr}=\dfrac{4\chi_mI}{R^2}$ тече вздовж $I$, повний $I'_\text{об}=4\pi\chi_mI$. Поверхневий
		$i_m=cJ(R)=\dfrac{2\chi_mI}{R}$ тече проти $I$, повний $I'_\text{пов}=4\pi\chi_mI$. Сумарний молекулярний струм крізь переріз дорівнює
		нулю: молекулярні струми замкнені й~заряду крізь переріз не переносять.
	\end{solution}
\end{problem}

\begin{problem}\label{prb:MagnetizedNeedlePlate}
	\emph{Однорідно намагнічене тіло.} (а)~Довгий тонкий стрижень (довжина $\gg$ діаметра) однорідно намагнічений вздовж осі з намагніченістю $J$. Знайдіть поверхневий молекулярний струм і поля $\Bfield$, $\Hfield$ усередині стрижня далеко від торців. (б)~Те саме для тонкої необмеженої пластини, однорідно намагніченої перпендикулярно до її площини. Підказка: скористайтеся неперервністю $H_\tau$ і $B_n$ на поверхні та тим, що поза тілом зовнішнього поля немає (у~(а) --- з точністю до малого поля торців).
	\begin{solution}
		(а)~$i_m=cJ$ (струм тече по колу навколо осі), $B=4\pi J$, $H=0$. (б)~Молекулярні струми течуть лише на краях пластини (нескінченно далеко), тому $B=0$, а $H=-4\pi J$.
	\end{solution}
\end{problem}

\begin{problem}\label{prb:MagneticCircuitGap}
	\emph{Магнітопровід із зазором.} Замкнений магнітопровід із магнетика з проникністю $\mu$ (довжина середньої лінії в магнетику $l$) має вузький повітряний зазор $\delta\ll l$; розсіянням поля й крайовими ефектами знехтуйте. На магнітопровід рівномірно намотано $N$ витків зі струмом $I$. Знайдіть $\Bfield$ і $\Hfield$ у магнетику та в зазорі. У скільки разів зменшується $B$ через зазор при $\mu=10^3$, $l=50$~см, $\delta=1$~мм?
	\begin{solution}
		З теореми про циркуляцію $H_\text{м}l+H_\text{з}\delta=\dfrac{4\pi}{c}NI$ і неперервності $B_n$ ($B=\mu H_\text{м}=H_\text{з}$) маємо $B=\dfrac{4\pi NI}{c\,(l/\mu+\delta)}$, $H_\text{м}=B/\mu$, $H_\text{з}=B$. Без зазору $B_0=\dfrac{4\pi\mu NI}{cl}$, тож $B_0/B=1+\dfrac{\mu\delta}{l}=3$.
	\end{solution}
\end{problem}

\begin{problem}\label{prb:ParaVsDia}
	\emph{Пара- чи діамагнетизм?} Речовина складається з атомів із концентрацією $n=5\cdot10^{22}$~см$^{-3}$, кожен з яких має один власний магнітний момент $p_0=\mu_B=9{,}3\cdot10^{-21}$~ерг/Гс і одному електрону відповідає $\langle r^2\rangle=10^{-16}$~см$^2$. Оцініть за порядком величини при $T=300$~К: (а)~парамагнітну сприйнятливість за законом Кюрі; (б)~діамагнітну сприйнятливість ($Z=1$); (в)~порівняйте їх; (г)~перевірте, що при $B=10^4$~Гс виконується $p_0B\ll kT$. ($k=1{,}38\cdot10^{-16}$~ерг/К, $e=4{,}8\cdot10^{-10}$~СГСЕ, $m_ec^2=8{,}2\cdot10^{-7}$~ерг.)
	\begin{solution}
		(а)~$\chi_m^{\text{пара}}=\dfrac{np_0^2}{3kT}\approx3{,}5\cdot10^{-5}$. (б)~$\chi_m^{\text{діа}}=-\dfrac{ne^2\langle r^2\rangle}{6m_ec^2}\approx-2{,}3\cdot10^{-7}$. (в)~Парамагнітний внесок більший приблизно в $1{,}5\cdot10^{2}$ разів, тож речовина парамагнітна. (г)~$p_0B/kT\approx2\cdot10^{-3}\ll1$.
	\end{solution}
\end{problem}

\begin{problem}\label{prb:WeissField}
	\emph{Молекулярне поле Вейса.} Вважаючи, що феромагнітне впорядкування руйнується, коли $p_0H_\text{мол}\sim kT_C$, оцініть поле Вейса для заліза ($T_C=1043$~К, $p_0\approx2{,}2\,\mu_B$, $\mu_B=9{,}3\cdot10^{-21}$~ерг/Гс) і порівняйте його з полем сусіднього атомного диполя на відстані $a=2{,}5\cdot10^{-8}$~см (за порядком величини $p_0/a^3$). Який висновок випливає щодо природи молекулярного поля?
	\begin{solution}
		$H_\text{мол}\sim\dfrac{kT_C}{p_0}\approx7\cdot10^{6}$~Гс, тоді як $\dfrac{p_0}{a^3}\approx1{,}3\cdot10^{3}$~Гс, тобто приблизно в $5\cdot10^{3}$ разів менше. Магнітної взаємодії диполів недостатньо для феромагнітного впорядкування; потрібна обмінна взаємодія.
	\end{solution}
\end{problem}

%\begin{problem}\label{prb:HysteresisLoss}
%	\emph{Втрати на гістерезис.} (а)~Робота джерела на перемагнічування одиниці об'єму феромагнетика за один цикл дорівнює $w=\dfrac1{4\pi}\oint H\,dB$. Покажіть, що $w=\oint H\,dJ$, тобто дорівнює площі петлі $J(H)$. (б)~Знайдіть $w$, замінивши петлю прямокутником з вершинами $(\pm H_c,\pm J_s)$, де $H_c=1$~Е (в гауссовій системі 1~Е чисельно відповідає 1~Гс), $J_s=1{,}6\cdot10^{3}$~Гс. (в)~Оцініть потужність втрат в осерді масою 1~кг ($\rho=7{,}8$~г/см$^3$), що перемагнічується з частотою 50~Гц.
%	\emph{Відповідь:} (а)~$B=H+4\pi J$, а $\oint H\,dH=0$ (повний диференціал по замкненому циклу). (б)~$w=4H_cJ_s=6{,}4\cdot10^{3}$~ерг/см$^3$. (в)~$V\approx128$~см$^3$, $P=w\nu V\approx4\cdot10^{7}$~ерг/с $\approx4$~Вт.
%\end{problem}


\begin{Questions}
\begin{quizlist}
	\item У чому полягала гіпотеза Ампера про природу магнетизму речовини? Як сучасна атомна фізика (рух електронів, їхній спін) пояснює, звідки в речовині беруться струми, про які здогадався Ампер?
	\item Що таке мікрополе $\Bfield_\text{micro}$ і середнє (макроскопічне) поле $\langle\Bfield\rangle$ у речовині? У чому полягає аналогія цієї конструкції з переходом від мікро- до макрополя в діелектриках?
	\item Чим струми провідності принципово відрізняються від молекулярних струмів?
	\item Дайте означення вектора намагніченості $\vect J$ як густини магнітних дипольних моментів. У чому полягає точна аналогія між $\vect J$ і вектором поляризації $\vect P$ в діелектриках?
	\item Чому пряме обчислення поля намагніченої речовини через підсумовування всіх молекулярних струмів практично нездійсненне, і яку ідею пропонує теорія магнетизму для розв'язання цієї проблеми?
	\item Виведіть формулу поверхневої густини молекулярного струму $i_m = c\,\vect J\cdot\vect\tau$, розглядаючи косокутний циліндр у намагніченій речовині. Чому молекулярний поверхневий струм визначається саме тангенціальною складовою $\vect J$, на відміну від зв'язаного заряду, що визначається нормальною складовою $\vect P$?
	\item Виведіть зв'язок об'ємної густини молекулярного струму з намагніченістю в диференціальній формі $\vect j_m = c\rot\vect J$, застосувавши теорему Стокса до циркуляції $\vect J$ по довільному контуру.
	\item Виведіть означення вектора напруженості магнітного поля $\Hfield = \Bfield - 4\pi\vect J$ з теореми про циркуляцію $\Bfield$, що враховує і струми провідності, і молекулярні струми. Запишіть теорему про циркуляцію в речовині в інтегральній та диференціальній формах.
	\item Чому вектор $\Hfield$ вважають допоміжною, а не фундаментальною величиною? У чому полягає аналогія між парами $(\Efield,\Dfield)$ і $(\Hfield,\Bfield)$ (таблиця~\ref{tab:DielMagAnalogy}) і чим вона неповна~--- дивергенція проти ротора?
	\item Запишіть лінійне співвідношення $\vect J = \chi_m\Hfield$ для лінійних ізотропних магнетиків та виведіть зв'язок $\Bfield=\mu\Hfield$, $\mu=1+4\pi\chi_m$.
	\item Чому термін \enquote{магнетик} не протиставляється нічому, на відміну від терміна \enquote{діелектрик}? Чим відрізняється наявність мікрострумів від наявності магнітного моменту атома?
	\item Дайте класифікацію лінійних магнетиків за знаком $\chi_m$: чим парамагнетики відрізняються від діамагнетиків за порядком величини сприйнятливості? Чим феромагнетики принципово відрізняються від обох цих класів?
	\item Виведіть першу граничну умову $B_{1n}=B_{2n}$ для нормальної складової $\Bfield$ на межі двох магнетиків із теореми Гаусса для $\Bfield$.
	\item Виведіть другу граничну умову $H_{2\tau}-H_{1\tau} = (4\pi/c)\,i_b$ для тангенціальної складової $\Hfield$ з теореми про циркуляцію. За яких умов $H_\tau$ неперервна, а $B_\tau$ --- ні?
	\item Виведіть закон заломлення силових ліній $\tg\alpha_2/\tg\alpha_1 = \mu_2/\mu_1$ на межі двох магнетиків.
	\item Що таке магнітне екранування і чим воно принципово відрізняється від електростатичного захисту? Чому феромагнітна оболонка \enquote{затягує} силові лінії в себе, а не виштовхує поле зі свого об'єму?
	\item Виведіть формулу ларморової частоти $\omega_L = eB/(2m_ec)$, розглядаючи прецесію орбітального моменту імпульсу електрона в зовнішньому полі. Чому ця частота однакова для будь-якого електрона незалежно від параметрів орбіти?
	\item Виведіть формулу індукованого магнітного моменту $\Delta p_m = -e^2\rho^2 B/(4m_ec^2)$, зумовленого прецесією Лармора. Чому цей момент завжди напрямлений проти поля, незалежно від початкового напрямку обертання електрона?
	\item Виведіть формулу Ланжевена для діамагнітної сприйнятливості $\chi_m^{\text{діа}} = -nZe^2\langle r^2\rangle/(6m_ec^2)$. Чому діамагнетизм --- універсальний ефект, присутній в атомах будь-якої речовини, і чому він не залежить від температури?
	\item Виведіть закон Кюрі для парамагнітної сприйнятливості $\chi_m^{\text{пара}} = C/T$, використовуючи розподіл Больцмана для орієнтації власних магнітних моментів атомів. Чому парамагнітний внесок, коли він є, зазвичай на порядки перевищує діамагнітний?
	\item У чому полягає гіпотеза Вейса про молекулярне поле $\Hfield_{\text{мол}} = \lambda\vect J$? Виведіть закон Кюрі--Вейса $\chi_m = C/(T-T_C)$ і поясніть, чому при $T\to T_C^+$ парамагнітний стан втрачає стійкість.
	\item Чому справжньою причиною феромагнітного впорядкування є обмінна, а не пряма магнітна взаємодія диполів? Як принцип Паулі пояснює залежність енергії електронів від взаємної орієнтації їхніх спінів?
	\item Чим відрізняються феро-, антиферо- та феримагнетики за характером обмінної взаємодії між сусідніми моментами?
	\item Опишіть петлю гістерезису феромагнетика: що таке намагніченість насичення $J_s$, залишкова намагніченість $J_r$ та коерцитивна сила $H_c$?
	\item Яким чином площа петлі гістерезису пов'язана з енергетичними втратами при перемагнічуванні? Чим магнітом'які матеріали відрізняються від магнітотвердих і де застосовується кожен з цих класів?
	\item Опишіть загальну логічну схему побудови теорії поля в речовині --- діелектриків і магнетиків: як вона проходить від густини мікроскопічних джерел до допоміжного вектора й лінійного співвідношення в обох випадках?
\end{quizlist}
\end{Questions}

\Closesolutionfile{answer}
